\section{Introduction}
\label{sec:intro}

Interstellar distances are measured in parsecs: one parsec is
$3.1\times10^{13}$~km, and even the nearest star, Proxima Centauri, is
\SunFirstHop~pc away. Any civilisation that spreads beyond its own system
therefore moves \emph{star to star}, and the natural mathematical object is a
graph. The stars are the nodes, a hop is an edge, and the engineering limits of
a ship (fuel, speed, how long the hardware or crew survives) cap the length of
the edges it can use.

This framing goes back to the Fermi-paradox literature
\cite{hart1975explanation,newman1981galactic}. Those studies asked how long a
fleet of probes would take to cover the Galaxy, on synthetic star fields with
the stars held fixed \cite{bjork2007exploring,cotta2009computational}. Two things
have changed since then. First, Gaia now gives positions, proper motions and,
for most nearby stars, radial velocities \cite{gaia2023dr3,smart2021gcns}, so
the real 6D star field can be routed instead of a random one. Second,
Carroll-Nellenback et al.\ \cite{carroll2019fermi} argued that stellar motions
alone can carry settlement across the Galaxy. That claim was made with a
statistical front model, not with the routes themselves.

We ask four concrete questions on the real catalogue.
\begin{enumerate}
\item \textbf{Static reach.} With stars frozen, what hop range $r$ is needed
to reach a given fraction of the neighbourhood from the Sun, and where are the
bottlenecks?
\item \textbf{Moving stars.} If a colony may \emph{wait} for a neighbour to
drift within range, how much smaller can $r$ be? This is an exact
time-dependent shortest-path problem, and it is the main contribution.
\item \textbf{Engineering.} Which propulsion technologies, per-hop lifetimes,
refuelling rules and pilot heuristics turn that range into a route?
\item \textbf{Exploration.} How much do optimised multi-probe tours gain over
the greedy rule used by earlier models?
\end{enumerate}

\paragraph{Contributions.} (i) An exact reduction of minimum-range routing to
the bottleneck of the Euclidean MST, applied to \NStars\ Gaia stars, with
astrometric bootstrap errors, a Poisson null model and a test for unseen brown
dwarfs. (ii) A closed-form temporal graph for linearly moving stars, in which
every edge is open on exactly one time window, and an exact earliest-arrival
Dijkstra over it. Waiting for stars lowers the percolation range by up to
\RedThousand\% within 1~Myr. (iii) A propulsion-to-range map, an exact
fuel-bottleneck algorithm for refuelling only at exoplanet hosts, and a
luminosity-limited photon-sail network. (iv) A reproduction of the classic
probe-exploration model on real stars. Every number in the text is generated
from the code and checked by unit tests against brute-force references.
